\paragraph{C.3ai (Refuted, with explicit counterexamples): every nonempty low-dimensional ordinary restriction range.}
For every pair of integers
\[
       r\ge1,\qquad r+2\le N\le2r+1,
\]
there is a specified regular excellent strictly henselian local
complex ring $R$ of dimension $N$, a nonzero $f\in\mathfrak m_R$,
and regular punctures $X$ and $Y=X\cap V(f)$ for which
\[
    \mathrm{CH}^r(X)=0,
    \qquad \mathrm{CH}^r(Y)_{\mathrm{alg}}\otimes\mathbb Q\ne0.
\]
We prove this by explicit suspended Fermat hypersurfaces. In
particular we do not assume a generalized Hodge conjecture or
infer non-torsion from a finite computation.

\begin{proof}
Set
\[
 d=2r-N+2\ge1,\qquad k=N-r-2\ge0,
 \qquad e=d+2,
 \qquad r=d+k,
 \qquad N=d+2k+2.
\]
Over $k_0=\overline{\mathbb Q}\subset\mathbb C$ define
\[
 \begin{split}
 T&=V(y_0^e+\cdots+y_{d+1}^e)\subset\mathbb P^{d+1},\\
 V&=V\left(\sum_{j=0}^{k-1}(a_j^e+b_j^e)
                    +\sum_{i=0}^{d+1}y_i^e\right)
                         \subset\mathbb P^{d+2k+1}.
 \end{split}
\]
Thus $T$ is a Calabi--Yau $d$-fold, $V$ has dimension
$n=d+2k$, and the coordinate embedding $T\subset V$ sets
all $a_j,b_j$ equal to zero. Write $H=c_1(\mathcal O_V(1))$.
Choose $\eta,\zeta\in k_0$ such that $\eta^e=-1$ and
$\zeta$ is a primitive $e$th root of unity. The two distinct
scalars $\lambda_0=\eta$ and $\lambda_1=\eta\zeta$ both
have $e$th power $-1$.

For $k>0$, impose $a_j=\lambda_i b_j$ for every $j$.
The resulting subvariety $R_i\subset V$ is the join of $T$
with the linear vertex $\mathbb P^{k-1}$ in the paired
coordinates. Its incidence resolution $Z_i\subset T\times V$
is a $\mathbb P^k$-bundle over $T$: over a point
$[y_0:\cdots:y_{d+1}]$, its fibre consists of
\[
 [\lambda_i b_0:b_0:\cdots:\lambda_i b_{k-1}:b_{k-1}:
                         t y_0:\cdots:t y_{d+1}].
\]
These coordinates describe the projective bundle homogeneously.
The fibre lies on $V$ since $\lambda_i^e+1=0$.
Its dimension is $d+k=r$, so we obtain
\[
       \Gamma_i=[Z_i]\in\mathrm{CH}^{r}(T\times V).
\]
For $k=0$ we put $V=T$ and
$\Gamma_0=\Gamma_1=\Delta_T$.

We have an integral left inverse
\begin{equation}\label{eq:C3-suspended-left-inverse}
             {}^t\Gamma_1\circ\Gamma_0=\Delta_T
                         \quad\text{in }\mathrm{CH}^{d}(T\times T).
\end{equation}
For $k>0$, the two different linear equations in each pair
force $a_j=b_j=0$. Thus $R_0\cap R_1=T$ scheme-theoretically,
and the intersection avoids both vertices. At a point of
$T$, the tangent space of $V$ is the direct sum of $T_T$
and the $2k$ paired-coordinate directions. The tangent
spaces of the two joins use complementary $k$-dimensional
subspaces of those directions, since
$\lambda_0\ne\lambda_1$. Their intersection is therefore
transverse. The fibre product of the two incidence
resolutions over $V$ is exactly the diagonal copy of $T$,
with multiplicity one. The Chow composition formula proves
\eqref{eq:C3-suspended-left-inverse}. The assertion is immediate
for $k=0$.

Choose $u_1,\ldots,u_d$ algebraically independent over $k_0$,
choose $w$ with $w^e=-1-\sum u_i^e$, and put
\[
    P=(u_1:\cdots:u_d:w:1),\qquad
    o=(1:\eta:0:\cdots:0),\qquad
    \alpha=[P]-[o]\in\mathrm{CH}^{d}(T)_{\mathrm{alg}}.
\]
The point $P$ is a geometric generic point over $k_0$.
The ordinary algebraic triviality is supplied by an actual
smooth complete-intersection curve through $P,o$ when
$d>1$, and by $T$ itself when $d=1$. Thus
$(\Gamma_0)_*\alpha\in\mathrm{CH}^{r}(V)_{\mathrm{alg}}$.
We claim that its image in
\[
       \mathrm{CH}^{r}(V)_{\mathbb Q}/
                          H\mathrm{CH}^{r-1}(V)_{\mathbb Q}
\]
is nonzero.

Suppose otherwise. Descend the finitely many cycles and
rational functions in a relation
$(\Gamma_0)_*\alpha=H\beta$ to the algebraic closure of
$k_0(T)$, then to a finite extension of that field, and
take the norm. Here the descent from $\mathbb C$ uses
specialization of finite data over the algebraically closed
subfield, as in C.3s. Closing the resulting generic-fibre
cycles on $T\times V$ and applying localization gives
\begin{equation}\label{eq:C3-suspended-spread}
 b\Gamma_0
    =H B+b\bigl(T\times(\Gamma_0)_*[o]\bigr)+\Theta,
 \qquad b\ne0,
\end{equation}
where $B\in\mathrm{CH}^{r-1}(T\times V)_{\mathbb Q}$ and
$\Theta\in\mathrm{CH}^{r}(T\times V)_{\mathbb Q}$ is
supported on $Z\times V$ for a proper divisor $Z\subset T$.
We retain the constant correspondence term explicitly;
it is not absorbed into the support condition on $\Theta$.

Apply these correspondences to $H^{d,0}(T)\ne0$. The
constant correspondence on the right acts as zero, since
$d>0$. The term supported on $Z\times V$ also acts as
zero: restriction to a resolution of $Z$ kills a
holomorphic $d$-form. The codimension of $B$ is $r-1$,
so its action sends this space into
\[
       H^{d+k-1,k-1}(V)
                   \subset H^{d+2k-2}(V,\mathbb C).
\]
For $k=0$ this is zero because its second Hodge index is
negative. For $k\ge1$ the total degree is $n-2<n$.
Weak Lefschetz identifies it with the same cohomology of
projective space, which has only diagonal Hodge types.
For the precise range, the complement of $V$ in
$\mathbb P^{n+1}$ is affine. The vanishing in C.2b and
Poincare duality give $H_c^j(\mathbb P^{n+1}\setminus V)=0$
for $j<n+1$. The compact-support localization sequence
therefore makes restriction an isomorphism in every degree
less than $n$, compatibly with Hodge structures.
The two displayed Hodge indices differ by $d\ge1$, so
this space is again zero. It follows that $HB$ acts as
zero on $H^{d,0}(T)$.

On the other hand, \eqref{eq:C3-suspended-left-inverse}
makes $(\Gamma_0)_*$ injective on cohomology; its left
inverse is $({}^t\Gamma_1)_*$. Thus the left side of
\eqref{eq:C3-suspended-spread} cannot act as zero on the
nonzero space $H^{d,0}(T)$. This contradiction proves
the claim. The argument uses weak Lefschetz below the
middle degree and the explicit correspondence identity;
no assertion about algebraic classes in the middle is needed.

Finally specify the regular local ambient ring and its
equation by
\[
 \begin{split}
 R&=\mathbb C[a_0,b_0,\ldots,a_{k-1},b_{k-1},
                y_0,\ldots,y_{d+1}]_{(a,b,y)}^{\mathrm{sh}},\\
 f&=\sum_{j=0}^{k-1}(a_j^e+b_j^e)+\sum_{i=0}^{d+1}y_i^e.
 \end{split}
\]
The paired-coordinate list is omitted when $k=0$.
The ring has dimension $d+2k+2=N$. The hypersurface is normal
and regular off its vertex, by the diagonal Fermat
gradient and the hypersurface depth criterion. Thus
$Y=(\operatorname{Spec}R\cap V(f))\setminus\{\mathfrak m_R\}$
is the strict-local punctured cone on $V$ and is regular.
The projectively normal cone injection C.3e sends the
nonzero quotient class just proved to a non-torsion class
in $\mathrm{CH}^{r}(Y)_{\mathrm{alg}}$. Ordinary algebraic
triviality follows by pulling the actual curve family and
the displayed incidence correspondence to the fixed
puncture. Corollary~6.4(i) of \citet[author p.~27]{rootSkalit2017},
with base $\mathbb C$, gives
$\mathrm{CH}^r(\operatorname{Spec}R)=0$; localization then
gives $\mathrm{CH}^r(X)=0$. This proves the assertion for
every pair in the stated range.
\end{proof}

The endpoints $k=0$ and $d=1$ recover respectively the
Calabi--Yau construction C.3ah and the odd-dimensional
Fermat correspondence in C.3r. The full interval between
them is proved by the same cohomological degree calculation.

The retained program is
\begin{center}
\texttt{code/check\_suspended\_calabi\_yau.py}.
\end{center}
It checks the parameter conversions and equations, as well as the
tangent matrices, cycle codimensions and Hodge shifts.
The all-dimensional non-torsion argument is the proof above.
